\section{导读：这期为什么是 AI for Math 的入口课}

这期访谈主角洪乐潼 Carina Hong，是 Axiom Math 创始人，方向是 AI for Math：让 AI 参与数学问题求解、证明形式化、证明搜索和数学研究协作。视频标题里的关键词很多：Lean、数学天书中的证明、直觉、被创造与被发现、Axiom、Ken Ono、Neo Lab。这些词共同指向一个问题：如果 AI 要进入真正数学研究，它要学会的不只是解题，还要理解数学语言、证明结构、形式系统和数学家的直觉。

\begin{importantbox}{本期核心命题}
AI for Math 的关键，不是让模型做更多竞赛题，而是把数学直觉、非形式证明、形式化语言、proof assistant 和机器验证接成闭环。Lean 不是装饰，而是把“看起来对”变成“机器可检查”的关键接口。
\end{importantbox}

\begin{knowledgebox}{视觉策略说明}
本视频是固定访谈画面，没有教学 slides、白板或产品演示。按本仓库播客标准，正文不重复插入人物帧；封面用于来源识别，正文用概念图、路线图和术语表承载教学内容。
\end{knowledgebox}

\subsection{本章小结}

这期不是普通创业访谈，而是 AI for Math 的概念入口：它同时讨论数学哲学、形式化证明、Lean、创业路线、研究直觉和 Neo Lab 的资源结构。


